\subsection{Schroeder--Bernstein and Cardinal Comparison}

\begin{theorem} \label{thm:schroeder-bernstein} \hlt{Schroeder--Bernstein Theorem}\\
If there are injections $f:A\to B$ and $g:B\to A$, then there is a bijection $h:A\to B$.
\begin{proof}
Let $A_0\defeq A\setminus g(B)$, and recursively define $B_n\defeq f(A_n)$ and $A_{n+1}\defeq g(B_n)$.\\
The sets $A_n$ are pairwise disjoint: if $A_i\cap A_j\neq\emptyset$ with $i<j$, then $i=0$ already contradicts $A_j\subseteq g(B)$; if $i>0$, injectivity of $g$ and then $f$ steps backward to $A_{i-1}\cap A_{j-1}\neq\emptyset$, and iteration gives the same contradiction.\\
Since $f$ is injective, the sets $B_n=f(A_n)$ are pairwise disjoint as well.\\
Let $A_*\defeq A\setminus\bigcup\limits_{n\geq0}A_n$ and $B_*\defeq B\setminus\bigcup\limits_{n\geq0}B_n$. The map $g$ restricts to a bijection $B_*\to A_*$.\\
Indeed, $g(B_*)$ avoids every $A_n$, while if $a\in A_*$ then $a\notin A_0$, so $a=g(b)$ for a unique $b\in B$; this $b$ cannot lie in any $B_n$, otherwise $a\in A_{n+1}$.\\
Define $h(x)\defeq
\begin{cases}
f(x), & x\in\bigcup_{n\geq0}A_n,\\
g^{-1}(x), & x\in A_*,
\end{cases}$, where $g^{-1}$ denotes the inverse of $g:B_*\to A_*$.\\
This maps each $A_n$ bijectively onto $B_n$ and maps $A_*$ bijectively onto $B_*$. Hence $h$ is a bijection $A\to B$.\\
\hlt{Structural Commentary}: Two injections carve both sets into matched layers plus a common remainder.\\
\hlt{Key Insight}: Use $f$ on the escaping layers and $g^{-1}$ on the stable remainder.
\end{proof}
\end{theorem}

\begin{definition} \label{def:cardinal-comparison-and-arithmetic} \hlt{Cardinal Comparison and Arithmetic}\\
Let $A,B$ be sets, and let $\kappa=\abs{A}$, $\lambda=\abs{B}$.
\begin{enumerate}[label=(\roman*)]
\setlength{\itemsep}{0pt}
\item \hlt{Cardinal comparison}: $\abs{A}\leq\abs{B}$ iff $\exists$ an injection $A\to B$.
\item \hlt{Strict cardinal comparison}: $\abs{A}<\abs{B}$ iff $\abs{A}\leq\abs{B}$ and $A\not\sim_{\#}B$.
\item If $A\cap B=\emptyset$, the \hlt{cardinal sum} is $\kappa+\lambda\defeq\abs{A\cup B}$.
\item The \hlt{cardinal product} is $\kappa\lambda\defeq\abs{A\times B}$.
\item The \hlt{function set} is $A^B\defeq\{f:B\to A\}$, and the \hlt{cardinal power} is $\kappa^\lambda\defeq\abs{A^B}$.
\end{enumerate}
\end{definition}
